\documentclass[twoside,openright]{report}


\usepackage[english]{babel}
\usepackage{kotex}    

\usepackage[letterpaper,top=1.2cm,bottom=1.2cm,left=1.2cm,right=1.2cm,marginparwidth=1.75cm]{geometry}

% Useful packages
\usepackage{amsmath, amsthm, amssymb, graphicx, enumitem, authblk, tikz, tikz-cd, verbatim, relsize}
\usepackage{thmtools}
\usepackage{tcolorbox}
\usepackage{forest} % 트리 구조를 위한 패키지
\usepackage{mathtools}
\usepackage[colorlinks=true, allcolors=black]{hyperref}
\usepackage[framemethod=TikZ]{mdframed}
\usepackage{lmodern}

\DeclareMathAlphabet{\mathtt}{OT1}{lmtt}{m}{n}

\usepackage{fontspec}
\setmainfont{neodgm_code.ttf}[
  Path = ./,
  UprightFont = *,
  Ligatures = TeX,
  Scale = MatchLowercase,
  % 폰트에 Bold/Italic이 없을 때 임시로:
  AutoFakeBold = 2.0,
  AutoFakeSlant = 0.2
]

\usetikzlibrary{calc}
\usetikzlibrary{shapes,backgrounds}
\usetikzlibrary{decorations.pathreplacing}
\usetikzlibrary{arrows.meta, positioning, matrix, fit}

% \usetikzlibrary{arrows.meta, positioning, external}
% \pgfplotsset{compat=1.18} % Set compatibility for pgfplots

\newtheoremstyle{romanstyle} % 스타일 이름
  {10pt} % 위쪽 여백
  {10pt} % 아래쪽 여백
  {\rmfamily} % 본문 로마체
  {} % 들여쓰기 없음
  {\bfseries} % 제목 굵게
  {.} % 제목 뒤 구두점
  { } % 제목과 본문 사이의 기본 간격
  {} % 제목 레이아웃 기본값

% 새 스타일 적용
\theoremstyle{romanstyle}


% 1) 박스형 스타일 프리셋
\declaretheoremstyle[
  headfont=\bfseries, bodyfont=\rmfamily,
  headpunct={.}, postheadspace=.5em,
  mdframed={linewidth=0.6pt, linecolor=black, backgroundcolor=white,
            innerleftmargin=8pt, innerrightmargin=8pt,
            innertopmargin=6pt, innerbottommargin=6pt,
            skipabove=10pt, skipbelow=10pt, roundcorner=0pt}
]{thm-plain}

\declaretheoremstyle[
  headfont=\bfseries, bodyfont=\rmfamily, headpunct={.},
  postheadspace=.5em,
  mdframed={linewidth=1pt, linecolor=black, backgroundcolor=gray!4,
            leftline=true, rightline=false, topline=false, bottomline=false,
            innerleftmargin=10pt, innerrightmargin=8pt,
            innertopmargin=6pt, innerbottommargin=6pt,
            skipabove=10pt, skipbelow=10pt}
]{thm-left}

\declaretheoremstyle[
  headfont=\bfseries, bodyfont=\rmfamily, headpunct={.},
  postheadspace=.5em,
  mdframed={linewidth=0.8pt, linecolor=black, backgroundcolor=white,
            roundcorner=2pt, tikzsetting={dash pattern=on 3pt off 2pt},
            innerleftmargin=8pt, innerrightmargin=8pt,
            innertopmargin=6pt, innerbottommargin=6pt}
]{thm-dashed}

\declaretheoremstyle[
  headfont=\bfseries, bodyfont=\rmfamily, headpunct={.},
  postheadspace=.5em,
  mdframed={linewidth=0.8pt, linecolor=blue!60!black, backgroundcolor=black!3,
            roundcorner=1.5pt,
            innerleftmargin=8pt, innerrightmargin=8pt,
            innertopmargin=6pt, innerbottommargin=6pt}
]{thm-blue}

\declaretheorem[style=thm-left,   numberwithin=subsection]{theorem}
\declaretheorem[style=thm-blue,   numberwithin=subsection]{proposition}
\declaretheorem[style=thm-dashed, numberwithin=subsection]{lemma}
\declaretheorem[style=thm-blue,   numberwithin=subsection]{corollary}
\declaretheorem[style=thm-plain,  numberwithin=subsection]{definition}

\newtheorem{axiom}{Axiom}

% Unnumbered environments
\newtheorem*{observation}{Observation}
\newtheorem*{example}{Example}
\newtheorem*{solution}{Solution}

% Mathematical commands
\newcommand{\st}{ \ \operatorname{s.t} \ }
\newcommand{\as}{ \ \operatorname{as} \ }
\newcommand{\setx}{\{x_0, x_1, \dots, x_n\}}
\newcommand{\setn}{\{1, \dots, n\}}
\newcommand{\dis}{\displaystyle}
\newcommand{\lete}{Let $\epsilon > 0$ be given.}
\newcommand{\fs}{\text{ for some }}
\newcommand{\defequal}{\overset{\textnormal{def}}{=}}
\newcommand{\setequal}{\overset{\textnormal{set}}{=}}
\newcommand*{\aand}{_{\text{ and }}}
\newcommand*{\oor}{_{\text{ or }}}
\DeclareMathOperator{\argmax}{argmax}
\DeclareMathOperator{\argmin}{argmin}
\newcommand{\im}{\operatorname{im}}
\newcommand{\nullity}{\operatorname{nullity}}
\newcommand{\rank}{\operatorname{rank}}
\newcommand{\ipvs}{\langle \, \cdot \, \rangle}
\newcommand{\diam}{\operatorname{diam}} % Add this line in the preamble
\newlength{\datedwd}
\newcommand{\dated}[2]{%
  \settowidth{\datedwd}{#1}%
  \noindent
  \parbox[t]{\datedwd}{\strut #1}%
  \hspace{0.5em}%
  \begin{minipage}[t]{\dimexpr\linewidth-\datedwd-0.5em\relax}
    \vspace{-0.65em}%
    \begin{enumerate}[label=\arabic*., leftmargin=*, itemsep=0pt, topsep=0pt, parsep=0pt]
      #2
    \end{enumerate}
  \end{minipage}\par
}

% \addparttoc{제목} 형태로 사용하도록 정의
\newcommand{\addparttoc}[1]{%
    \phantomsection
    \addtocontents{toc}{\protect\addvspace{1.0\baselineskip}} % 목차 내 여백 추가
    \addcontentsline{toc}{part}{#1}% 목차에 'part' 레벨로 추가
}


% ===================== TOC에 정리 제목(옵션 있을 때만) 자동 추가 ===================== %
\usepackage{xparse}   %% <<< ADDED
\usepackage{etoolbox} %% <<< ADDED
\setcounter{tocdepth}{3}    %% <<< ADDED: subsubsection까지 TOC 표시
\setcounter{secnumdepth}{3} %% <<< ADDED: (선택) 번호도 3단계까지

% theorem
\let\theoremOLD\theorem           %% <<< ADDED
\let\endtheoremOLD\endtheorem     %% <<< ADDED
\RenewDocumentEnvironment{theorem}{o}{%% <<< ADDED
  \IfNoValueTF{#1}{%
    \theoremOLD
  }{%
    \theoremOLD
    \phantomsection
    \addcontentsline{toc}{subsubsection}{%
      \protect\numberline{\thetheorem}\ #1}%
  }%
}{\endtheoremOLD}

% lemma
\let\lemmaOLD\lemma
\let\endlemmaOLD\endlemma
\RenewDocumentEnvironment{lemma}{o}{%
  \IfNoValueTF{#1}{%
    \lemmaOLD
  }{%
    \lemmaOLD
    \phantomsection
    \addcontentsline{toc}{subsubsection}{%
      \protect\numberline{\thelemma}\ #1}%
  }%
}{\endlemmaOLD}

% proposition
\let\propositionOLD\proposition
\let\endpropositionOLD\endproposition
\RenewDocumentEnvironment{proposition}{o}{%
  \IfNoValueTF{#1}{%
    \propositionOLD
  }{%
    \propositionOLD
    \phantomsection
    \addcontentsline{toc}{subsubsection}{%
      \protect\numberline{\theproposition}\ #1}%
  }%
}{\endpropositionOLD}

% corollary
\let\corollaryOLD\corollary
\let\endcorollaryOLD\endcorollary
\RenewDocumentEnvironment{corollary}{o}{%
  \IfNoValueTF{#1}{%
    \corollaryOLD
  }{%
    \corollaryOLD
    \phantomsection
    \addcontentsline{toc}{subsubsection}{%
      \protect\numberline{\thecorollary}\ #1}%
  }%
}{\endcorollaryOLD}

% definition
\let\definitionOLD\definition
\let\enddefinitionOLD\enddefinition
\RenewDocumentEnvironment{definition}{o}{%
  \IfNoValueTF{#1}{%
    \definitionOLD
  }{%
    \definitionOLD
    \phantomsection
    \addcontentsline{toc}{subsubsection}{%
      \protect\numberline{\thedefinition}Definition\ #1}%
  }%
}{\enddefinitionOLD}
% ==================================================================== %


\title{Math Note}
\author{Jong Won}
\affil{University of Seoul, Mathematics}
\date{}
% \sloppy
\begin{document}
\maketitle

\newpage
\tableofcontents

\newpage
\noindent Patch Note: \\
\dated{2026/03/03 -}{
\item End of TEPS. Start Mathematics.
}



\addparttoc{Algebra}
\newpage

\chapter{Prelininary}
\begin{definition}
    Let $G$ be a set. \\
    A map $\cdot: G \times G \to G : (a, b) \mapsto a \cdot b$ is called \textit{Binary Operation} on $G$. \\
    The Binary Operation is said to be:
    \begin{enumerate}
        \item \textit{Associative} if: for all $a, b, c \in G$, $(a \cdot b) \cdot c = a \cdot (b \cdot c)$.
        \item \textit{Commutative} if: for all $a, b \in G$, $a \cdot b = b \cdot a$.
        \item An element $e \in G$ is called \textit{identity} if: for all $g \in G,\ g \cdot e = e \cdot g = g$.
        \item An element $g \in G$ has \textit{inverse} if: for all $g \in G$, there exists $g^{-1} \st g^{-1} \cdot g = e\ \operatorname{or}\ g \cdot g^{-1} = e$.
        \item Binary Operation $\times:G \times G \to G$ is \textit{distributive} with respect to a Binary Operation $+:G \times G \to G$ if:
        \begin{align*}
            \text{For all $a,b,c \in G$, }
            \begin{cases}
                (a + b) \times c = a \times c + b \times c \\
                a \times (b + c) = a \times b + a \times c
            \end{cases}
        \end{align*}
    \end{enumerate}
\end{definition}

\begin{definition}
    Let $G, A$ be sets. \\
    A map $\ast: G \times A \to A:(g, a) \mapsto g \ast a$ is called \textit{Action} of $G$ on $A$. \\
    The Action is said to be with respect to a Binary Operation $+, \cdot:G \times G \to G$ if:
    \begin{enumerate}
        \item Associative with respect to $\cdot$ if: for all $g_1, g_2 \in G, a \in A,\ (g_1 \cdot g_2) \ast a = g_1 \ast (g_2 \ast a)$.
        \item Distributive with respect to $+$ if: for all $g_1, g_2 \in G, a \in A,\ (g_1 + g_2) \ast a = g_1 \ast a +_A g_2 \ast a$.
        \item Distributive with respect to $+_A$ if: for all $g \in G, a_1, a_2 \in A,\ g \ast (a_1 +_A a_2) = g \ast a_1 +_A g \ast a_2$.
    \end{enumerate}
\end{definition}

\chapter{Group Theory}
\begin{definition}
    Let $G$ be a set, and $\cdot:G \times G \to G:(a,b) \mapsto a \cdot b$ be a Binary Operation on $G$. \\
    A Tuple $(G, \ \cdot\ )$ is called \textit{Group} if: 
    \begin{enumerate}
        \item The Opearation $\cdot$ is Associative.
        \item $G$ has an identity.
        \item Every element of $G$ has an inverse.
    \end{enumerate}
\end{definition}
\section{General Properties}

\section{Cyclic Group}

\section{Normal Subgroup}

\section{Isomorphism Theorems}

\section{Group Action}
\begin{definition}
    Let $(G, \ \cdot \ )$ be a Group, and $A$ be a set. \\
    An action $\ast:G \times A \to A$ is called \textit{Group Action} if:
    \begin{enumerate}
        \item For all $g_1, g_2 \in G,\ a \in A,\ g_1 \ast(g_2 \ast a) = (g_1 \cdot g_2) \ast a$.
        \item For any $a \in A,\ 1 \ast a = a$.
    \end{enumerate}
\end{definition}
\section{Automorphism}

\section{Commutator Subgroup}

\section{Free Group}

\chapter{Finite Group Theory}

\section{Class Equation}

\section{Sylow Theorem}

\section{$p$-Group}

\subsection{Existence of $p$-Subgroup}

\section{Hölder program}




\chapter{Ring Theory}

\section{Ideal}

\section{Ring of Fractions}
\begin{theorem}
    Let $R$ be a Commutative Ring, $D \subset R$ be a subset such that $\begin{cases}
        \text{no zero, no zero divisors} \\ \text{closed under multiplication}
    \end{cases}$. \\
    Then, there exists a Commutative Ring $Q$ with identity satisfies:
    \begin{enumerate}
        \item $R$ can embed in $Q$, and every element of $D$ becomes unit in $Q$.
        More precisely, $Q = \{r d^{-1} \mid r \in R,\ d \in D\}$.
        \item $Q$ is the smallest Ring contaning $R$ with identity such that every element of $D$ becomes unit in $Q$.
    \end{enumerate}
\end{theorem}
\begin{proof}
    Let $\mathcal{F} \defequal \{(r,d) \mid r \in R,\ d \in D\}$ and the relation $\sim$ on $\mathcal{F}$
    by $(r_1, d_1) \sim (r_2, d_2) \iff r_1d_2 = r_2d_1$.\\
    Then, $\sim$ is equivalent relation: reflexive and symmetirc are clear, and
    Suppose that $(r_1, d_1) \sim (r_2, d_2)$ and $(r_2, d_2) \sim (r_3, d_3)$.
    \begin{align*}
        r_2d_3 = r_3d_2 \implies r_2 d_1 d_3 = r_3 d_1 d_2 \implies r_1 d_2 d_3 = r_3 d_1 d_2 \implies d_2 ( r_1 d_3 - r_3 d_1) \implies r_1d_3 = r_3d_1
    \end{align*}
    Thus transitivity shown. Define
    \begin{align*}
        \frac{r}{d} \defequal [(r,d)] = \{(a,b) \mid (a,b) \sim (r,d) \},\ Q \defequal \left\{ \frac{r}{d} \;\middle|\; r \in R,\ d \in D \right\}
    \end{align*}
    And define operations $+, \times$ on $Q$:
    \begin{align*}
        \frac{r_1}{d_1} + \frac{r_2}{d_2} \defequal \frac{r_1d_2 + r_2d_1}{d_1 d_2}, \quad
        \frac{r_1}{d_1} \times \frac{r_2}{d_2} \defequal \frac{r_1r_2}{d_1d_2}
    \end{align*}
    Well-Definedness: If $\dis \frac{r_1}{d_1} = \frac{r'_1}{d'_1}$ and $\dis \frac{r_2}{d_2} = \frac{r'_2}{d'_2}$,
    \begin{align*}
        \frac{r_1d_2 + r_2d_1}{d_1 d_2}
        = \frac{r_1 d_2 d'_1 d'_2 + r_2 d_1 d'_1 d'_2}{d_1 d_2 d'_1 d'_2}
        = \frac{(r_1 d'_1) d_2 d'_2 + (r_2 d'_2) d_1 d'_1 }{d_1 d_2 d'_1 d'_2}
        = \frac{(r'_1 d_1) d_2 d'_2 + (r'_2 d_2) d_1 d'_1 }{d_1 d_2 d'_1 d'_2}
        = \frac{(r'_1d'_2 + r'_2d'_1)d_1d_2}{d_1d_2 d'_1 d'_2}
        = \frac{r'_1d'_2 + r'_2d'_1}{d'_1 d'_2}
    \end{align*}
    \begin{align*}
        \frac{r_1r_2}{d_1d_2} 
        = \frac{r_1 r_2 d'_1 d'_2}{d_1d_2d'_1d'_2}
        = \frac{(r_1 d'_1) (r_2 d'_2)}{d_1d_2d'_1d'_2}
        = \frac{(r'_1 d_1) (r'_2 d_2)}{d_1d_2d'_1d'_2}
        = \frac{r'_1 r'_2 d_1 d_2}{d_1d_2d'_1d'_2}
        = \frac{r'_1 r'_2}{d'_1d'_2}
    \end{align*}
    Now, $(Q, +, \times)$ constructs Commutative Ring with identity: for any $d \in D$, 
    put $\dis 0_Q \defequal \frac{0}{d},\ 1_Q \defequal \frac{d}{d}$. Then,
    \begin{enumerate}
        \item $(R, +, \times)$ closed under the operations since $D$ is closed under the multiplication.
        \item $(R, +)$ has a zero: $\dis \frac{r_1}{d_1} + 0_Q = \frac{r_1}{d_1} + \frac{0}{d} = \frac{r_1 d + 0 d_1}{d_1 d} = \frac{r_1 d}{d_1 d} = \frac{r_1}{d_1}$.
        \item $(R, +)$ has an inverse: $\dis \frac{r_1}{d_1} + \frac{-r_1}{d_1} = \frac{r_1d_1 + (-r_1)d_1}{d_1 d_1} = \frac{[(r_1) + (-r_1)]d_1}{d_1 d_1} = \frac{0 d_1}{d_1 d_1} = \frac{0}{d_1 d_1} = 0_Q$.
        \item $(R, +, \times)$ satisfies distributive law:
        \begin{enumerate}[label=\theenumi-\arabic*.]
            \item The left law:
            \begin{align*}
                \frac{r_1}{d_1} \times \left( \frac{r_2}{d_2} + \frac{r_3}{d_3} \right)
            = &\frac{r_1}{d_1} \times \frac{r_2 d_3 + r_3 d_2}{d_2 d_3}
            = \frac{r_1 r_2 d_3 + r_1 r_3 d_2}{d_1 d_2 d_3}
            = \frac{r_1 r_2 d_1 d_3 + r_1 r_3 d_1 d_2}{d_1 d_2 d_1 d_3}
            = \frac{r_1 r_2}{d_1 d_2} + \frac{r_2 r_3}{d_2 d_3} \\
            = &\frac{r_1}{d_1} \times \frac{r_2}{d_2} + \frac{r_2}{d_2} \times \frac{r_3}{d_3}
            \end{align*}
            \item The right law:
            \begin{align*}
                \left(\frac{r_1}{d_1} + \frac{r_2}{d_2}\right) \times \frac{r_3}{d_3}
            = &\frac{r_1 d_2 + r_2 d_1}{d_1 d_2} \times \frac{r_3}{d_3}
            = \frac{r_1 r_3 d_2 + r_2 r_3 d_1}{d_1 d_2 d_3}
            = \frac{r_1 r_3 d_2 d_3 + r_2 r_3 d_1 d_3}{d_1 d_3 d_2 d_3}
            = \frac{r_1 r_3}{d_1 d_3} + \frac{r_2 r_3}{d_2 d_3} \\
            = &\frac{r_1}{d_1} \times \frac{r_3}{d_3} + \frac{r_2}{d_2} \times \frac{r_3}{d_3}
            \end{align*}
        \end{enumerate}
        \item $(R, \times)$ has an identity: $\dis \frac{r_1}{d_1} \times 1_Q = \frac{r_1}{d_1} \times \frac{d}{d} = \frac{r_1d}{d_1d} = \frac{r_1}{d_1}$.
        \item Elements of $D$ become unit in $Q$: Define $\iota:R \to Q:r \mapsto \dis \frac{rp}{p}$ where $p \in D$ is any fixed element in $D$. \\
        Then, $\iota$ is Ring-Monomorphsim because:
        \begin{enumerate}[label=\theenumi-\arabic*.]
            \item Well-Defined and Injective: 
            $\dis \iota (r_1) = \iota (r_2)
            \iff \frac{r_1 p}{p} = \frac{r_2 f}{f}
            \iff (r_1 - r_2) pp = 0
            \iff r_1 = r_2$
            \item For any $d \in D$, $\iota(d)$ is a unit of $Q$: Put $\dis (\iota(d))^{-1} \defequal \frac{p}{dp}$, then
            \begin{align*}
                \iota(d) \times (\iota(d))^{-1} = \frac{dp}{p} \times \frac{p}{dp} = \frac{dpp}{dpp} = 1_Q
            \end{align*}
            \noindent That is, $\iota$ is embedding from $R$ into $Q$ such that $\iota[D]$ becomes units of $Q$ except zero. \\
            Moreover, if $D = R \setminus \{0\}$, then $Q$ is field.
        \end{enumerate}
        \item $Q$ is the \textit{smallest} ring containing $R$ with identity such that every element of $D$ becomes units in $Q$. \\
        Let $S$ be an any commutative ring with identity, \\
        and assume that $\varphi:R \to S$ is a Ring-Monomorphism such that for any $d \in D$, $\varphi(d)$ is unit in $S$. \\
        Define $\dis \phi:Q \to S:\frac{r}{d} \mapsto \varphi(r) \varphi(d)^{-1}$. Then, this $\phi$ is well-defined and injective:
        \begin{align*}
            &\phi \left(\frac{r_1}{d_1} \right) = \phi \left(\frac{r_2}{d_2} \right)
            \iff
            \varphi(r_1) \varphi(d_1)^{-1} = \varphi(r_2) \varphi(d_2)^{-1}
            \iff 
            \varphi(r_1) \varphi(d_2) = \varphi(r_2) \varphi(d_1)\\
            \overset{\text{homomor.}}{\iff}
            &\varphi(r_1 d_2) = \varphi(r_2 d_1) 
            \overset{\text{one-to-one}}{\iff}
            r_1 d_2 = r_2 d_1
            \iff
            \frac{r_1}{d_1} = \frac{r_2}{d_2}
        \end{align*}
        That is, if a commutative ring $S$ with identity contains a copy of $R$ such that the denominator set $D$ of $R$ becomes unit in $S$, then $S$ contains ring of fractions $Q$ of $R$.
        Thus $S = Q$ is the smallest ring that satisfies these conditions.
    \end{enumerate}
\end{proof}

\section{Commutative Ring with Unity}

\chapter{Polynomial Ring Theory}

\chapter{Module Theory}

\section{Basic Definitions}

\begin{definition}
    Let $R$ be a Ring. A \textit{left Module over $R$} is an Abelain Group $(M, +)$ with: \\
        An action $\cdot$ of $R$ on $M$ satisfies: for any $r,s \in R,\ m,n \in M$,
        $\begin{cases}
            (r + s) \cdot m = r \cdot m + s \cdot m \\
            (r s) \cdot m = r \cdot (s \cdot m) \\
            r \cdot (m + n) = r \cdot m + r \cdot n 
        \end{cases}$. \\
    In addition, if $R$ has an identity $1$, $1 \cdot m = m$
\end{definition}

\begin{definition}
    Let $M$ be a left Module over a Ring $R$. \\
    A subgroup $N \leq M$ is called \textit{submodule} of $M$ if: which is closed under the action of $R$ on $M$. \\
    That is, for all $r \in R$, $n \in N$, $r \cdot n \in N$.
\end{definition}

\begin{definition}
    Let $M, N$ be lef Modules over a Ring $R$. \\
    A map $\varphi:M \to N$ is called \textit{Module-Homomorphism} if: for any $x,y \in M, r \in R$,
    $\begin{cases}
        \varphi(x + y) = \varphi(x) + \varphi(y) \\
        \varphi(r \cdot x) = r \cdot \varphi(x)
    \end{cases}$.
\end{definition}

\begin{lemma}
    \textbf{Submodule Criteria} \\
    Let $M$ be a left Module over a Ring $R$. A subgroup $N \leq M$ is a submodule of $M$ if and only if
    \begin{align*}
        N \neq \emptyset,\
        \forall x, y \in N,\ r \in R,\ r \cdot x + y \in N
    \end{align*}
\end{lemma}


\begin{definition}
    Let $R$ be a Ring with identity. \\
    A Ring $A$ with identity is called \textit{$R$-Algebra} if: there exists a Ring Homomorphism $f:R \to A$ such that
    \begin{align*}
        f(1_R) = 1_A,\ Z(A) \subseteq f[R]
    \end{align*}
    Remark: $Z(A) \defequal \{r \in A \mid \forall a \in A,\ ar = ra\}$.
\end{definition}

\begin{definition}
    Let $A, B$ be $R$-Algebras. \\
    A Ring-Homomorphism $\varphi:A \to B$ is called \textit{Algebra-Homomorphism} if: $\varphi(1_A) = 1_B$ and
    \begin{align*}
        \forall r \in R, a \in A,\ \varphi(r \cdot a) = r \cdot \varphi(a)
    \end{align*}
\end{definition}

\newpage
\section{Module of Module-Homomorphism}
\begin{definition}
    Let $M,N$ be left Modules over a Ring $R$. \\
    Define two operation on $\operatorname{Hom}_R(M,N)$: for any $\varphi, \psi \in \operatorname{Hom}_R(M,N), r \in R$,
    \begin{align*}
        \begin{cases}
            \varphi + \psi:M \to N:m \mapsto \varphi(m) + \psi(n) \\
            r \cdot \varphi: M \to N : m \mapsto r \cdot \varphi(m)
        \end{cases}
    \end{align*}
    Then, $(\operatorname{Hom}_R(M,N), +, \ \cdot\ , R)$ becomes a left Module over $R$.
\end{definition}

\newpage
\section{Direct Sum of Modules}
\begin{definition}
    Let $M$ be a left Module over a Ring $R$ and let $N_1, \ldots, N_n$ be submodules of $M$. \\
    Define \textit{Sum} of $N_i$:
    \begin{align*}
        N_1 + N_2 + \cdots + N_n \defequal \{a_1 + a_2 + \cdots + a_n \mid a_i \in N_i,\ \forall 1 \leq i \leq n \}
    \end{align*}
    Let $A$ be any subset of $M$. Define \textit{submodule generated by $A$}:
    \begin{align*}
        RA \defequal \{r_1 \cdot a_1 + \cdots + r_n \cdot a_n \mid n \in \mathbb{N}, r_i \in R, a_i \in A\}
    \end{align*}
    A submodule $N \subseteq M$ is said to be \textit{finitely generated}: for some finite subset $A \subset M,\ N = RA$. \\
    A submodule $N \subseteq M$ is called \textit{cyclic} if: there exists $a \in M$ such that $N = R a = \{r \cdot a \mid r \in R\}$.
\end{definition}


\chapter{Vector Space Theory}
\section{Vector Space}
\begin{definition}
    Suppose that $F$ is a Field, and $V$ is an Abelian Group. \\
    And, the operation $\cdot: F \times V \to V$ satisfies: For any $a, b \in F$ and $v, w \in V$,
    $\begin{cases}
        a \cdot (v + w) = a\cdot v + a\cdot w \\
        (a + b) \cdot v = a\cdot v + b\cdot v \\
        (ab) \cdot v = a \cdot (b \cdot v) \\
        1 \cdot v = v
    \end{cases}$ \\
    The triple $(V, +, \cdot)$ is called the \textit{Vector Space} over $F$.
\end{definition}
Equivalently, The Vector Space over a field $F$ is $F$-Module.
\begin{definition}
    Suppose that $V$ is a Vector Space over a field $F$, and $W \subseteq V$ is a Subset. \\
    The $W$ is called \textit{Subspace} of $V$ if:
    $\begin{cases}
        \text{For any $a \in F,\ w \in W,\ a \cdot w \in W$} \\
        \text{For any $w_1, w_2 \in W,\ w_1 + w_2 \in W$}
    \end{cases}$. \\
\end{definition}
That is, the Subset of a Vector Space which is closed under the addition and scalar multiplication, then it is a Vector Space.
\begin{lemma}
    Arbitrary intersection of Subspace is a Subspace.
\end{lemma}
\begin{proof}
    Suppose that $V$ is a Vector Space, and $W_\alpha \leq V,\ \alpha \in \Lambda$ are Subspaces. \\
    Using Subspace Criterion: Let $a \in F,\ w_1, w_2 \in \bigcap_{\alpha \in \Lambda} W_\alpha$ be given.
    Then, for any $\alpha \in \Lambda$, $a \cdot w_1 + w_2 \in W_\alpha$.
\end{proof}
This Lemma allows the definition:
\begin{definition}
    Suppose that $V$ is a Vector Space over a field $F$, and $S \subseteq V$ be a Subset. \\
    Define a \textit{Generated Subspace} by $S$ is:
    \begin{align*}
        \langle S \rangle \defequal \bigcap_{S \subseteq W \leq V} \hspace{-0.2cm} W
    \end{align*}
\end{definition}
This $\langle S \rangle$ is the \textit{unique smallest} Subspace containing $S$. This $S$ is called \textit{Generating Subset} of $\langle S \rangle$.
\begin{lemma}
    Suppose that $V$ is a Vector Space over a field $F$, and $S \subseteq V$ be a Subset. Then,
    \begin{align*}
        \langle S \rangle = \{ a_1 \cdot v_1 + \cdots + a_n \cdot v_n \mid n \in \mathbb{N},\ a_i \in F,\ v_i \in S \}
    \end{align*}
\end{lemma}
\begin{proof}
    First, $\langle S \rangle \supseteq \{ a_1 \cdot v_1 + \cdots + a_n \cdot v_n \mid n \in \mathbb{N},\ a_i \in F,\ v_i \in S \}$, because $\langle S \rangle$ is closed under the opeartions. \\
    And, the set $\{ a_1 \cdot v_1 + \cdots + a_n \cdot v_n \mid n \in \mathbb{N},\ a_i \in F,\ v_i \in S \}$ is a Subgroup of $V$ containing $S$, \\
    Hence $\langle S \rangle \subseteq \{ a_1 \cdot v_1 + \cdots + a_n \cdot v_n \mid n \in \mathbb{N},\ a_i \in F,\ v_i \in S \}$.
\end{proof}

\begin{definition}
    Let $V$ be a Vector Space over a field $F$, and $W \leq V$ be a Subspace. \\
    The \textit{Quotient Space} $V/W$ is defined as the set of Cosets of $W$ in $V$:
    \begin{align*}
        V/W \defequal \{v + W \mid v \in V\}
    \end{align*}
    with the operations: for any $v, v_1, v_2 \in V,\ a \in F$,
    \begin{align*}
        (v_1 + W) + (v_2 + W) &\defequal (v_1 + v_2) + W \\
        a \cdot (v + W) &\defequal (a \cdot v) + W
    \end{align*}
\end{definition}
The well-definedness of above definition is allowed by the follwings. \\
Observation: 
for any $v, v' \in V$,
\begin{align*}
    v + W = v' + W 
    \iff v - v' \in W
    \iff \exists w \in W \st v - v' = w
\end{align*}
Now, suppose that $v_1, v_2, v'_1, v'_2 \in V$ 
with $v_1 + W = v'_1 + W$ and $v_2 + W = v'_2 + W$. \\
Then, there exist $w_1, w_2 \in W$ such that $v_1 - v'_1 = w_1$ and $v_2 - v'_2 = w_2$. Since above observation,
\begin{align*}
    (v_1 + W) + (v_2 + W)
    = ((v'_1 + w_1) + W) + ((v'_2 + w_2) + W)
    = (v'_1 + W) + (v'_2 + W)
\end{align*}
And, similrarly, the scalr-multiplication is well-defined.

\begin{lemma}
    Let $V$ be a Vector Space over $F$, and $W_1, W_2 \leq $ be subspaces. Then,
    \begin{enumerate}
        \item $W_1 + W_2$ is the smallest subspace of $V$ that contains both $W_1$ and $W_2$.
    \end{enumerate}
\end{lemma}
\begin{proof}
    It is clear that $W_1 + W_2$ is a subspace of $V$ by the Subspace Criterion. \\
    To show that it is the smallest such space, Suppose that $W \leq V$ is a subspace that contains both $W_1$ and $W_2$. \\
    Then, for any $w_1 + w_2 \in W_1 + W_2$ with $w_i \in W_i$, $w_1 + w_2 \in W$ because closed under the operation.
\end{proof}

\newpage

\begin{definition}
    Suppose that $V$ is a Vector Space over $F$, and $S \subseteq V$ is a subset. \\
    A vector $v \in V$ is called \textit{linear combination of $S$} if: 
    there exist finite numbers $a_i \in F$ and $v_i \in S$ such that
    \begin{align*}
        v = a_1 v_1 + a_2 v_2 + \cdots + a_n v_n
    \end{align*}
\end{definition}

\section{Linearly independent}
\begin{definition}
    Suppose that $V$ is a Vector Space over a field $F$, and $S \subseteq V$ is a Subset. \\
    A Subset $S$ is called \textit{Linearly independent} if: For any finite subset $\{v_1, v_2, \ldots, v_n\} \subseteq S$,
    \begin{align*}
        a_1 \cdot v_1 + a_2 \cdot v_2 + \cdots + a_n \cdot v_n = 0 \implies a_1 = a_2 = \cdots = a_n = 0
    \end{align*}
    If $S$ is not Linearly independent, then it is called \textit{Linearly dependent}.
\end{definition}

% \begin{lemma}
%     Suppose that $S \subseteq V$ is a Linearly independent subset of a Vector Space $V$ over $F$. \\
%     If $v \in V$ is linear combination of $S$ with for some $a_i \in F,\ v_i \in S$,
%     \begin{align*}
%         v = a_1 v_1 + a_2 v_2 + \cdots + a_n v_n
%     \end{align*}
%     then this representation is unique. Briefly, linear combination of a vector is unique. 
% \end{lemma}
% \begin{proof}
%     First, $v = a_1 \cdot v_1 + \cdots + a_n \cdot v_n$ implies $v \in \langle S \rangle$. \\
%     Now, suppose that $v \in V$ satisfies $v = a_1 \cdot v_1 + \cdots + a_n \cdot v_n = b_1 \cdot w_1 + \cdots + b_m \cdot w_m$, WLOG $n \leq m$. \\
%     Put $I \defequal \{i \in \mathbb{N} \mid \exists j \in \mathbb{N} \st v_i = w_j\}$. \\
%     Then, there is a permutation $\rho: \{1, 2, \ldots, m\} \to \{1, 2, \ldots, m\}$ such that $\forall i \in I,\ v_i = w_{\rho(i)}$. Now,
%     \begin{align*}
%         &\sum_{i \in I} a_i \cdot v_i + \sum_{j \notin I} a_j \cdot v_j
%         = \sum_{i \in I} b_{\rho(i)} \cdot w_{\rho(i)} + \sum_{j \notin I} b_{\rho(j)} \cdot w_{\rho(j)}
%         = \sum_{i \in I} b_{\rho(i)} \cdot v_i + \sum_{j \notin I} b_{\rho(j)} \cdot w_{\rho(j)} \\
%         \implies &\sum_{i \in I} (a_i - b_{\rho(i)}) \cdot v_i + \sum_{j \notin I} a_j \cdot v_j - \left(\sum_{j \notin I} b_{\rho(j)} \cdot w_{\rho(j)}\right) = 0
%     \end{align*}
%     Since $S$ is linearly independent, for all $j \notin I,\ a_j = b_{\rho(j)} = 0$ and for all $i \in I,\ a_i = b_{\rho(i)}$.
%     % 1) If $\{v_1, \ldots, v_n \} \cap \{w_1, \ldots, w_m\} = \emptyset$. \\
%     % \begin{align*}
%     %     &a_1 \cdot v_1 + \cdots + a_n \cdot v_n = b_1 \cdot w_1 + \cdots + b_m \cdot w_m \\
%     %     \iff &a_1 \cdot v_1 + \cdots + a_n \cdot v_n - b_1 \cdot w_1 - \cdots - b_m \cdot w_m = 0 \\
%     %     \implies &a_1 = \cdots = a_n = b_1 = \cdots = b_m = 0
%     % \end{align*}
%     % This means $v \in V$ is a zero vector. \\
%     % 2) If $\{v_1, \ldots, v_n \} \cap \{w_1, \ldots, w_m\} \neq \emptyset$. \\
%     % Put $\{u_i \mid 1 \leq i \leq j\} = \{v_1, \ldots, v_n \} \cap \{w_1, \ldots, w_m\}$ where $j \leq n$. \\
%     % \begin{align*}
%     %     &a_{1} \cdot u_1 + \cdots + a_{j} \cdot u_j + a_{j+1} \cdot v_{j+1} + \cdots + a_n \cdot v_n
%     %     = b_{1} \cdot u_1 + \cdots + b_{j} \cdot u_j + b_{j+1} \cdot w_{j+1} + \cdots + b_m \cdot w_m \\
%     %     \iff &(a_{1} - b_{1}) \cdot u_1 + \cdots + (a_{j} - b_{j}) + a_{j+1} \cdot v_{j+1} + \cdots + a_n \cdot v_n 
%     %     - b_{j+1} \cdot w_{j+1} - \cdots - b_m \cdot w_m = 0 \\
%     %     \implies & a_{1} = b_{1},\ \ldots,\ a_{j} = b_{j},\ a_{j+1} = \cdots = a_n = b_{j+1} = \cdots = b_m = 0
%     % \end{align*}
%     % This means $v = a_1 u_1 + \cdots + a_j u_j$ is the Unique representation of $v \in V$ by linear combination of $S$.
% \end{proof}
% This fact enables the definition in the next section.

\newpage
\section{Basis}
\begin{definition}
    Suppose that $V$ is a Vector Space over a field $F$. \\
    A subset $\beta \subseteq V$ is called the \textit{Basis} of $V$ if:
    \begin{enumerate}
        \item $\beta$ is Linearly independent. 
        \item $\langle \beta \rangle = V$.
    \end{enumerate}
\end{definition}
\begin{lemma}
    Suppose that $V$ is a Vector Space over a field $F$. Then,
    \begin{center}
        $\beta \subset V$ is a Basis of $V$
        $\iff$
        For any $v \in V$, there exists a Unique linear combination of $\beta$ that equals $v$.
    \end{center}
\end{lemma}
\begin{proof}
    Suppose that $\beta \subset V$ is a Basis of $V$. 
    Because $v \in V = \langle \beta \rangle$, $v$ can be represented a linear combination of $\beta$. \\
    This representation must be unique. Suppose that $v$ be represented in two different forms:
    \begin{align*}
        v = a_1 v_1 + \cdots + a_n v_n = b_1 v_1 + \cdots b_n v_n
    \end{align*}
    where $a_i, b_i \in F,\ v_i \in \beta$. Then,
    \begin{align*}
        (a_1 - b_1) v_1 + \cdots + (a_n - b_n) v_n = 0 \implies \forall 1 \leq i \leq n,\ a_i = b_i
    \end{align*}
    Convesely, suppose that For any $v \in V$, there uniquely exist $a_i \in F,\ v_i \in \beta$ such that
    \begin{align*}
        v = a_1 v_1 + \cdots + a_n v_n
    \end{align*}
    $\langle \beta \rangle = V$ is clear. And, if $\beta$ is linearly dependent, then
    there exist a representation of $\beta$ that equals zero:
    \begin{align*}
        0 = a_1 v_1 + \cdots + a_n v_n
    \end{align*}
    where $v_i \in \beta$, $a_i \in F$ not all zero.
    This is contradiction with uniqueness.
\end{proof}

\begin{theorem}
    Let $V$ be a Vector Space over $F$, and $W_1, W_2 \leq V$ be finite-dimensional subspaces. \\
    Then, followings are hold:
    \begin{enumerate}
        \item $\dim(W_1 \cap W_2) \leq \min\{\dim W_1, \dim W_2 \}$.
        \item $\dim(W_1 + W_2) \leq \dim W_1 + \dim W_2$.
    \end{enumerate}
\end{theorem}
\begin{proof}
    Proof of 2. Let $\beta_1, \beta_2$ be basis of $W_1, W_2$, respectively.
    Since $W_i = \langle \beta_i \rangle$,
    \begin{align*}
        W_1 + W_2 = \langle \beta_1 \rangle + \langle \beta_2 \rangle
        &= \{a_1 v_1 + \cdots + a_n v_n \mid n \in \mathbb{N}, v_i \in \beta_1, a_i \in F\}
        + \{b_1 w_1 + \cdots + b_m w_m \mid m \in \mathbb{N}, w_i \in \beta_2, b_i \in F\} \\
        &= \{c_1 u_1 + \cdots + c_k u_k \mid k \in \mathbb{N}, u_i \in \beta_1 \cup \beta_2, c_i \in F\}
        = \langle \beta_1 \cup \beta_2 \rangle
    \end{align*}
    Thus,
    \begin{align*}
        \dim(W_1 + W_2) = \sharp(\beta_1 \cup \beta_2) \leq \sharp \beta_1 + \sharp \beta_2 = \dim W_1 + \dim W_2
    \end{align*}
\end{proof}

\begin{theorem}
    Let $V$ be a Vector Space over $F$, and $W_1, W_2 \leq V$ be subspaces such that $V = W_1 \oplus W_2$. \\
    If $\beta_1, \beta_2$ is basis of $W_1, W_2$ respectively, then $\beta_1 \cap \beta_2 = \emptyset$ and $\beta_1 \cup \beta_2$ is basis of $V$.
\end{theorem}
\begin{proof}
    If $\beta_1 \cap \beta_2$ is not empty, the for some $x \in \beta_1 \cap \beta_2$, $x \in W_1 \cap W_2$. It is contradiction to $V = W_1 \oplus W_2$. \\
    And, $\beta_1 \cup \beta_2$ spans $V$ because:
    \begin{align*}
        \langle \beta_1 \cup \beta_2 \rangle = \langle \beta_1 \rangle + \langle \beta_2 \rangle = W_1 + W_2 = V
    \end{align*}
    Now, to show linearly independent, suppose that some linear combination of above vectors be zero. That is,
    \begin{align*}
        &a_1 v_1 + \cdots + a_n v_n = 0 \quad (a_i \in F, v_i \in \beta_1 \cup \beta_2)
    \end{align*}
    Then, all coefficients must be zero because $\beta_1, \beta_2$ are linearly independent.
\end{proof}
The below theorem is converse statement of above.
\begin{theorem}
    Let $V$ be a Vector Space over $F$, and $W_1, W_2$ be subspaces with disjoint basis $\beta_1, \beta_2$ respectively.
    If $\beta_1 \cup \beta_2$ is basis for $V$, then $V = W_1 \oplus W_2$. 
\end{theorem}
\begin{proof}
    Since $\beta_1 \cup \beta_2$ spans $V$, 
    \begin{align*}
        V = \langle \beta_1 \cup \beta_2 \rangle = \langle \beta_1 \rangle + \langle \beta_2 \rangle = W_1 + W_2
    \end{align*}
    And, to show $W_1 \cap W_2 = \{0\}$, suppose that there exists a non-zero vector $v \in W_1 \cap W_2$. \\
    Then, there exist two different representation of $v$ respect with to $\beta_1$ and $\beta_2$:
    \begin{align*}
        v = a_1 v_1 + \cdots + a_n v_n = b_1 w_1 + \cdots + b_m 
    \end{align*}
    where $a_i, b_i \in F, v_i \in \beta_1, w_i \in \beta_2$. But, the disjointness and linearly indenpendent implies all $a_i = b_i = 0$. \\
    Thus, $v = 0$. It is contradiction.
\end{proof}

\begin{theorem}
    Let $V$ be a Vector Space over $F$, $W \leq V$ be a subspace. Put $\beta, \gamma$ be basis of $V, W$ respectively. \\
    If $\gamma \subset \beta$, then the set of cosets
    \begin{align*}
        \mathcal{B} = \{v + W \mid v \in \beta \setminus \gamma\}
    \end{align*}
    is basis of a Quotient Space $V / W$.
\end{theorem}
\begin{proof}
    Denote $\beta = \{v_i \mid i \in I\}$ and $\gamma = \{v_i \mid i \in E\}$ where $E \subset I$.
\end{proof}

\newpage
\section{$\dagger$ Existence of Basis}
\begin{theorem}
    Every Vector Space has a Basis.
\end{theorem}
\begin{proof}
    Using Zorn's Lemma. Let $V$ be a Vector Space.
    Put $P \defequal \{S \subseteq V \mid \text{$S$ is linearly independent subset}\}$. \\
    Clearly, $\{v\} \in P$ for any $v \in V$. Thus, $(P, \subseteq)$ is non-emepty partially ordered set which has inclusion as order. \\
    Suppose $\mathcal{C} \subseteq P$ is a chain of $P$. If $\mathcal{C} = \emptyset$,
    then $\{v\}$ is upper bound of $\mathcal{C}$ for any $v \in V$. Thus, assume $\mathcal{C} \neq \emptyset$. \\
    To show that $\mathcal{C}$ has an upper bound in $P$, denote the union of all elements in the chain:
    \begin{align*}
        B = \bigcup_{\beta \in \mathcal{C}} \beta
    \end{align*}
    Clearly this is upper bound of $\mathcal{C}$. Hence, enough to show that $B$ is linearly independent; That is, $B \in P$. \\
    Suppose that $B$ is linearly dependent. Then, there exist finite vectors in $B$ and scalars such that
    \begin{align*}
        a_1 v_1 + \cdots + a_n v_n = 0
    \end{align*}
    for some $a_j$ is not zero.
    Since $\mathcal{C}$ is a Chain, there exists some subset of $\mathcal{C}$ which containing $\{v_1, \ldots, v_n \}$. \\
    Because for each $v_i$, there exists $\beta_i \in \mathcal{C}$ such that $v_i \in \beta_i$, and
    \begin{align*}
        \{v_1, \ldots, v_n \} \subseteq \bigcup_{i=1}^n \beta_i \overset{(*)}{\in} \mathcal{C}
    \end{align*}
    The inclusion $(*)$ is given by finite construction: for any $\beta_i, \beta_j \in \mathcal{C}$,
    $\beta_i \subseteq \beta_j$ or $\beta_i \supseteq \beta_j$. \\
    Thus, the union of two elements in the chain is either one. 
    This is contradiction that the finite union is linearly independent, being it is contained in $P$.
    In sum, $\mathcal{C}$ has an upper bound in $P$. \\
    By Zorn's Lemma,
    there exists a Maximal $\mathcal{M} \subsetneq V$ which is linearly independent, and satisfies
    \begin{center}
        If $\beta \subseteq V$ is linearly independent, then $\beta \subseteq \mathcal{M}$.
    \end{center}
    Lastly, if $\langle \mathcal{M} \rangle \subsetneq V$, then there exists $v \in V \setminus \langle \mathcal{M} \rangle$ such that
    $\mathcal{M} \subsetneq \mathcal{M} \cup \{v\}$. \\
    But, this implies $\mathcal{M} \cup \{v\}$ is linearly independent and
    proper supset of $\mathcal{M}$, contradiction.  \\
    Consquently, $\langle \mathcal{M} \rangle = V$.
\end{proof}

\section{Linear Map}
\begin{definition}
    Suppose that $V, W$ are Vector Space over a field $F$.\\
    A map $\mathcal{L}:V \to W$ is called \textit{Linear Map} if:
    \begin{center}
        For any $a \in F$ and $v_1, v_2 \in V$, $\mathcal{L}(a \cdot v_1 + v_2) = a \cdot \mathcal{L}(v_1) + \mathcal{L}(v_2)$
    \end{center}
\end{definition}

\begin{theorem}
    Let $V, W$ be Vector Space over $F$, and $\beta$ be a basis of $V$. \\
    For any map $f:\beta \to W$, there exists unique linear map $T:V \to W$ such that
    \begin{align*}
        \forall v \in \beta,\ T(v) = f(v)
    \end{align*}
    Specifically,
    \begin{align*}
        T:V \to W: \sum_{i=1}^{n} a_i \cdot v_i \mapsto \sum_{i=1}^{n} a_i \cdot f(v_i)
    \end{align*}
    where $\dis V = \left\{ \sum_{i=1}^{n} a_i \cdot v_i \;\middle|\; n \in \mathbb{N}, a_i \in F, v_i \in \beta \right\}$
\end{theorem}
\begin{proof}
    Enough to Show that: $T$ is linear, $T$ satisfies the Condition, and Uniqueness. \\
    Let $x,y \in V$ and $\alpha \in F$ be given. Then, the vectors $x,y$ can be represented unique linear combination forms:
    \begin{align*}
        x = \sum_{i=1}^{n} a_i \cdot v_i,\
        y = \sum_{i=1}^{n} b_i \cdot v_i
    \end{align*}
    where $a_i, b_i \in F,\ v_i \in \beta$. Now,
    \begin{align*}
        T(\alpha \cdot x + y)
        &= T\left( \alpha \cdot \left( \sum_{i=1}^{n} a_i \cdot v_i \right) + \left( \sum_{i=1}^{n} b_i \cdot v_i \right) \right)
        = T\left( \sum_{i=1}^{n} (\alpha a_i + b_i) \cdot v_i \right)  \\
        &= \sum_{i=1}^{n} (\alpha a_i + b_i) \cdot f(v_i)
        = \alpha \cdot \left(\sum_{i=1}^{n} a_i \cdot f(v_i) \right) + \left( \sum_{i=1}^{n} b_i \cdot f(v_i) \right)\\
        &= \alpha \cdot T\left(\sum_{i=1}^{n} a_i \cdot v_i\right) + T \left(\sum_{i=1}^{n} b_i \cdot v_i \right)
        = \alpha \cdot T(x) + T(y)
    \end{align*}
    Thus, $T$ is linear map. Since $\beta$ is basis, for all $v \in \beta,\ T(v) = f(v)$ clearly. \\
    And, to show the Uniqueness, suppose that $U:V \to W$ is a linear map such that $\forall v \in \beta,\ U(v) = f(v)$. \\
    Then, for any vector $x \in V$ which has a linear combination form $x = \sum_{i=1}^{n} a_i v_i$,
    \begin{align*}
        U(x) = U \left( \sum_{i=1}^{n} a_i v_i \right)
        = \sum_{i=1}^{n} a_i U(v_i)
        = \sum_{i=1}^{n} a_i f(v_i)
        = T(x)
    \end{align*}
    Consequently, $U = T$.
\end{proof}

\newpage
\begin{theorem}
    Let $V, W$ be Vector Space over $F$, and $T: V \to W$ be a linear map. 
    \\ Put $\beta \subset V$ be a basis of $V$. Then,
    \begin{align*}
        \im T = \langle T[\beta] \rangle
    \end{align*}
    Remark: in definition,
    \begin{align*}
        \im T \defequal \{T(v) \mid v \in V\},\
        \langle T[\beta] \rangle \defequal 
        \left\{a_1 T(v_1) + \cdots + a_n T(v_n) \;\middle|\; n \in \mathbb{N}, a_i \in F, v_i \in \beta \right\}
    \end{align*}
\end{theorem}
\begin{proof}
    Trivial.
\end{proof}

\section{Direct Sum and Projection}
\begin{definition}
    Let $V$ be a Vector Space over $F$, and $W_1, W_2 \leq V$ be subspaces. \\
    $V$ is called \textit{direct sum} of $W_1$ and $W_2$ if:
    $V = W_1 + W_2$ and $W_1 \cap W_2 = \{0\}$.
\end{definition}

\begin{theorem}
    Let $V$ be a Vector Space over $F$, and $W_1, W_2 \leq V$ be subspaces. Then, TFAE:
    \begin{align*}
        V = W_1 \oplus W_2 \iff
        \forall v \in V,\
        \exists ! x_1 \in W_1, x_2 \in W_2 \st
        v = x_1 + x_2
    \end{align*}
\end{theorem}
\begin{proof}
    Suppose that $V = W_1 \oplus W_2$. Enough to show: $v$ has unique representation. \\
    For some $x_1, y_1 \in W_1, x_2, y_2 \in W_2$, $v = x_1 + x_2 = y_1 + y_2$ implies
    $x_1 - y_1 = y_2 - x_2$. \\
    Now, $x_1 - y_1 \in W_1$ and $y_2 - x_2 \in W_2$ and $W_1 \cap W_2 = \{0\}$ implies $x_1 - y_1 = y_2 - x_2 = 0$. \\
    Convesely, enough to show that $W_1 \cap W_2 = \{0\}$. \\
    Suppose that there exists non-zero element $v \in W_1 \cap W_2$.
    Then, zero can be represented in two different forms:
    \begin{align*}
        0 = 0 + 0 = v + (-v)
    \end{align*}
    That is, the representation of zero by a sum of $W_1$ and $W_2$ is not unique, which is a contradiction.
\end{proof}

\begin{definition}
    Let $V$ be a Vector Space over $F$ and $W_1, W_2 \leq V$ be subspace such that $V = W_1 \oplus W_2$. \\
    The following map $T$ is called \textit{projection on $W_1$ along $W_2$}:
    \begin{align*}
        T:V \to V:x_1 + x_2 \mapsto x_1 \ \operatorname{where} \ x_1 \in W_1, x_2 \in W_2
    \end{align*}
    The well-definedness is allowed by above theorem.
\end{definition}
\begin{theorem}
    Let $T:V \to V$ be a projection on $W_1$ along $W_2$. Then,
    \begin{enumerate}
        \item $T$ is linear map.
        \item $W_1 = \{x \in V \mid T(x) = x\} = \im T$.
        \item $W_2 = \ker T$.
    \end{enumerate}
\end{theorem}
\begin{proof}
    Linearlity is clear: \\
    For any $x = x_1 + x_2 \in V$ and $y = y_1 + y_2 \in V$ where $x_1, y_1 \in W_1, x_2, y_2 \in W_2$, $\alpha \in F$,
    \begin{align*}
        T(\alpha x + y)
        = T(\alpha x_1 + \alpha x_2 + y_1 + y_2)
        = T(\alpha x_1 + y_1 + \alpha x_2 + y_2)
        = \alpha x_1 + y_1
        = \alpha T(x_1 + x_2) + T(y_1 + y_2)
        = \alpha T(x) + T(y)
    \end{align*}
    Next, $W_1 \subseteq \{x \in V \mid T(x) = x\}$ is clear.
    And, represent $x \in V$ as $x = x_1 + x_2$ where $x_i \in W_i$. \\
    Then, $T(x) = x$ implies $x_2 = 0$, thus $x \in W_1$. Therefore, $W_1 = \{x \in V \mid T(x) = x\}$. And,
    \begin{align*}
        \im T = T[V] 
        = T[W_1 \oplus W_2] 
        = \{T(x_1 + x_2) \mid x_1 \in W_1, x_2 \in W_2\}
        = W_1
    \end{align*}
    $W_2 \subseteq \ker T$ is clear. Suppose that there exists non-zero $v \in W_2$ with $v \notin \ker T$. \\
    Then, $T(v) = 0$ and $v \notin \ker T$, which is contradiction.
\end{proof}

\begin{theorem}
    Let $V$ be a finite-dimensional Vector Space over $F$, and $W \leq V$ be a subspace. \\
    Then, there exists a subspace $W'$ and a projection $T:V \to V$ on $W$ along $W'$.
\end{theorem}
\begin{proof}
    Let basis of $V$ be $\beta = \{v_i \mid 1 \leq i \leq n\}$.
\end{proof}

%Concentrate to General Theory of Vector Space and Linear map.

\chapter{Field Theory}

\chapter{Galois Theory}

\addparttoc{Number Theory}
\chapter{Number Theory}
Well-Ordering Principle: Every non-empty subset of non-negative integers set has a minimum element.
\section{Definition}
\begin{definition}
    Let $a, b \in \mathbb{Z}$ be integers. \\
    We said to be $b$ divides $a$ if: $a = bx$ for some $x \in \mathbb{Z}$. Denote $b \mid a$.
\end{definition}

\section{The Division Algorithm}
\begin{theorem}
    Let $a, b \in \mathbb{Z}$ be integers with $b > 0$. 
    Then, there exists unique integers $q, r \in \mathbb{Z}$ such that
    \begin{align*}
        a = b \cdot q + r \quad (0 \leq r < b)
    \end{align*}
\end{theorem}
\begin{proof}
    Let $S \defequal \{a - bx \mid x \in \mathbb{Z},\ a - bx \geq 0\}$.
    Then, $S$ is non-empty because $a - b \cdot (-|a|) = a + |a| b \geq a + |a| \geq 0$. \\
    Thus, Well-Ordering Principle gives that there exists a minimum element of $S$. Put $r = \min S$. \\
    (That is, for any $k \in S$, $r \leq k$.)
    Since $r \in S$, there exists $q \in \mathbb{Z}$ such that $ r = a - bq \geq 0$. \\
    To show $r < b$, we will use Contradiction. Suppose that $r \geq b$. Then,
    \begin{align*}
        a - b \cdot (q + 1) = a - bq - b = r - b \geq 0
    \end{align*}
    Thus, $a - b (q + 1) \in S$, but this is contradiction with $r = a - bq$ is a minimum element of $S$. \\
    To show the Uniqueness, there exist $q, q', r, r' \in \mathbb{Z}$ such that
    \begin{align*}
        a = bq + r = bq' + r' \quad(0 \leq r, r' < b)
    \end{align*}
    Then, $b|q - q'| = |r - r'|$, thus must be $q = q'$ and $r = r'$.
\end{proof}
\begin{corollary}
    Let $a, b \in \mathbb{Z}$ be integers with $b \neq 0$. 
    Then, there exists unique integers $q, r \in \mathbb{Z}$ such that
    \begin{align*}
        a = b \cdot q + r \quad (0 \leq r < |b|)
    \end{align*}
\end{corollary}

\newpage
\section{Bézout's Identity}
\begin{definition}
    Let $a,b \in \mathbb{Z}$ be integers. The integer $d$ is called \textit{The greatset common divisor} of $a, b$ if:
    \begin{enumerate}
        \item $d \mid a$ and $d \mid b$.
        \item If $c \mid a$ and $c \mid b$, then $c \leq d$.
    \end{enumerate}
    Denote $d \defequal \gcd(a,b)$.
\end{definition}

\begin{theorem}
    Let $a,b \in \mathbb{Z}$ be integers. Then, there exist integers $x, y \in \mathbb{Z}$ such that
    \begin{align*}
        \gcd(a,b) = ax + by
    \end{align*}
\end{theorem}
\begin{proof}
    Define a set:
    \begin{align*}
        S \defequal \{a \cdot u + b \cdot v \mid u,v \in \mathbb{Z},\ a \cdot u + b \cdot v > 0\}
    \end{align*}
    Clearly, $S$ is non-emptyset. Thus, $S$ has a minimum element $d = \min S$ by Well-Ordering Principle. \\
    Since $d \in S$, $d = ax + by$ for some $x,y \in \mathbb{Z}$. By the Division Algorithm, there exist $q, r \in \mathbb{Z}$ such that
    \begin{align*}
        a = d \cdot q + r \quad(0 \leq r < d)
    \end{align*}
    That is,
    \begin{align*}
        r = a - d \cdot q = a - q \cdot (ax + by) = a \cdot (1 - qx) + b \cdot (-qy)
    \end{align*}
    If $r \neq 0$, then $r \in S$ and $r < d$, this is contradiction with $d = \min S$. Thus, $r = 0$. \\
    Similarly, $b = d \cdot q'$ for some $q' \in \mathbb{Z}$. In summary, $d$ divides $a$ and $b$. \\
    Now, let $c$ be an arbitrary positive common divisor of $a$ and $b$.
    Then, $c$ divides $ax + by = d$, thus $c \leq d$. \\
    Consequently, $d = \gcd(a,b)$.
\end{proof}
\begin{corollary}
    Let $a,b$ be integers. For $d = \gcd(a,b)$,
    \begin{align*}
        \{a \cdot x + b \cdot y \mid x, y \in \mathbb{Z}\} = \{d \cdot k \mid k \in \mathbb{Z}\}
    \end{align*}
\end{corollary}

\newpage
\subsection{Extended Euclidean Algorithm}
\begin{lemma}
    Let $a,b$ be integers. \\
    If $a = b \cdot q + r$ with $0 \leq r < b$,
    then $\gcd(a,b) = \gcd(b,r)$.
\end{lemma}
\begin{proof}
    Put $d = \gcd(a,b)$. Since $d \mid a$ and $d \mid b$, $d$ divides $r = a - bq$.
    That is, $d$ is common divisor of $b$ and $r$. \\
    Meanwhile, let $c$ be a divisor of $b$ and $r$. Then, $c$ divides $a = bq + r$, thus $c \leq d$.
    This implies $d = \gcd(b,r)$.
\end{proof}
Consider the given integers $a,b$. By the Division Algorithm, we can write:
\begin{align*}
    a &= b \cdot q_0 + r_0 &&(0 \leq r_0 < b)\\
    b &= r_0 \cdot q_1 + r_1 &&(0 \leq r_1 < r_0)\\
    r_0 &= r_1 \cdot q_2 + r_2 &&(0 \leq r_2 < r_1)\\
    &\ \vdots \\
    r_{k-2}&= r_{k-1} \cdot q_{k} + r_{k} &&(0 \leq r_k < r_{k-1})\\
    &\ \vdots \\
    r_{n-2} &= r_{n-1} \cdot q_n + r_n &&(0 \leq r_n < r_{n-1})\\
    r_{n-1} &= r_{n} \cdot q_{n+1} &&(r_{n+1} = 0)
\end{align*}
Since the Non-negative numbers Sequnece $\{r_k\}$ strictly decrease, $r_{n+1}$ must be zero for some $n \in \mathbb{N}$. \\
Put $r_{-2} = a$ and $r_{-1} = b$. Since above Lemma,
\begin{align*}
    \gcd(a,b) = \gcd(b, r_0) = \gcd(r_0, r_1) = \cdots = \gcd(r_{n-1}, r_n) = \gcd(r_n, 0) = r_n
\end{align*}
This implies also for any $-2 \leq k \leq n$, $r_k$ is multiple of $r_n$. \\
% By Bézout’s Identity, there exists $x, y \in \mathbb{Z}$ such that $r_n = \gcd(a,b) = a \cdot x + b \cdot y$. \\
% Meanwhile, for $k = n$, $r_{n-1} = r_n \cdot q_{n+1}$ is multiple of $r_n$.
% Put $r_{-2} = a$ and $r_{-1} = b$. Suppose that for $0 \leq k \leq n$, $r_{k-1}$ and $r_k$ both are multiple of $r_{n}$. Then, the formular
% \begin{align*}
%     r_{k-2} = r_{k-1} \cdot q_n + r_k
% \end{align*}
% implies $r_{k-2}$ is multiple of $r_n$. And, for $k = n$, $r_{n-1} = r_n \cdot q_{n+1}$ is multiple of $r_n$. \\
% Now, the Induction implies for all $0 \leq k \leq n$, $r_k$ is multiple of $r_n$. \\
By Corollary of Bézout’s Identity, for any $-2 \leq k \leq n$, there exists $x_k, y_k \in \mathbb{Z}$ such that
\begin{align*}
    r_k = a \cdot x_k + b \cdot y_k \tag{$*$}
\end{align*}
Since $r_{-2} = a$ and $r_{-1} = b$,
\begin{align*}
    &x_{-2} = 1,\ y_{-2} = 0 \\
    &x_{-1} = 0,\ y_{-1} = 1
\end{align*}
And, substituting $(*)$ to $r_{k-2} - r_{k-1} \cdot q_k = r_k$:
\begin{align*}
    a \cdot x_{k} + b \cdot y_{k}
    = a \cdot x_{k-2} + b \cdot y_{k-2} - (a \cdot x_{k-1} + b \cdot y_{k-1}) \cdot q_{k}
    = a \cdot (x_{k-2} - x_{k-1} \cdot q_k) + b \cdot (y_{k-2} - y_{k-1} \cdot q_k)
\end{align*}
Now, we obtain the equation: for all $0 \leq k \leq n$,
\begin{align*}
    &x_k = x_{k-2} - x_{k-1} \cdot q_k \\
    &y_k = y_{k-2} - y_{k-1} \cdot q_k
\end{align*}
Consequently, to find $x,y \in \mathbb{Z}$ such that $\gcd(a,b) = r_n = a \cdot x + b \cdot y = a \cdot x_n + b \cdot y_n$, \\
we simply calculate $x_n, y_n$ inductively.

% \begin{proof}
    % By the Division Algorithm, there exist $q_0, r_0$ such that $a = b \cdot q_0 + r_0$ with $0 \leq r_0 < b$. \\
    % Using again Division Algorithm for $b$ and $r_0$, there exist $q_1, r_1$ such that $b = r_0 \cdot q_1 + r_1$ with $0 \leq r_1 < r_0$. 
    % Inductively,
    % \begin{align*}
    %     a &= b \cdot q_0 + r_0 \\
    %     b &= r_0 \cdot q_1 + r_1 \\
    %     r_0 &= r_1 \cdot q_2 + r_2 \\
    %     & \vdots \\
    %     r_{k-2}&= r_{k-1} \cdot q_{k} + r_{k} \\
    %     & \vdots \\
    %     r_{n-2} &= r_{n-1} \cdot q_n + r_n \\
    %     r_{n-1} &= r_{n} \cdot q_{n+1}
    % \end{align*}
    % Since the Non-negative numbers Sequnece $\{r_k\}$ strictly decrease, $r_{n+1}$ must be zero for some $n \in \mathbb{N}$. \\
    % And, since above Lemma,
    % \begin{align*}
    %     \gcd(a,b) = \gcd(b, r_0) = \gcd(r_0, r_1) = \cdots = \gcd(r_{n-1}, r_n) = \gcd(r_n, 0) = r_n
    % \end{align*}
    % Now, rewrite the above formulas:
    % \begin{align*}
    %     r_n &= r_{n-2} - r_{n-1} \cdot q_n\\
    %     &= r_{n-2} - (r_{n-3} - r_{n-2} \cdot q_{n-1}) = r_{n-2} \cdot (1 + q_{n-1}) - r_{n-3} \\
    %     &= (r_{n-4} - r_{n-3} \cdot q_{n-2}) \cdot (1 + q_{n-1}) - r_{n-3}
    %     = r_{n-4} \cdot (1 + q_{n-1}) - r_{n-3} \cdot (1 + q_{n-2} + q_{n-1} \cdot q_{n-2}) \\
    %     &= r_{n-4} \cdot (1 + q_{n-1}) - (r_{n-5} - r_{n-4} \cdot q_{n-3}) \cdot (1 + q_{n-2} + q_{n-1} \cdot q_{n-2}) \\
    %     &= r_{n-4} \cdot (1 + q_{n-1} + q_{n-3} + q_{n-2} \cdot q_{n-3} + q_{n-1} \cdot q_{n-2} \cdot q_{n-3}) - r_{n-5} \cdot (1 + q_{n-2} + q_{n-1} \cdot q_{n-2}) \\
    %     &= (r_{n-6} - r_{n-5} \cdot q_{n-4}) \cdot (1 + q_{n-1} + q_{n-3} + q_{n-2} \cdot q_{n-3} + q_{n-1} \cdot q_{n-2} \cdot q_{n-3}) - r_{n-5} \cdot (1 + q_{n-2} + q_{n-1} \cdot q_{n-2}) \\
    %     &= r_{n-6} \cdot (1 + q_{n-1} + q_{n-3} + q_{n-2} \cdot q_{n-3} + q_{n-1} \cdot q_{n-2} \cdot q_{n-3})
    %     - r_{n-5} \cdot (1 + q_{n-2} + q_{n-1} \cdot q_{n-2} )
    % \end{align*}
    
    % \begin{align*}
    %     a - b \cdot q_0 &= r_0 \\
    %     b - r_0 \cdot q_1 &= r_1 \\
    %     r_0 - r_1 \cdot q_2 &= r_2 \\
    %     & \vdots \\
    %     r_k - r_{k+1} \cdot q_{k+2} &= r_{k+2} \\
    %     & \vdots \\
    %     r_{n-4} - r_{n-3} \cdot q_{n-2} &= r_{n-2} \\
    %     r_{n-3} - r_{n-2} \cdot q_{n-1} &= r_{n-1} \\
    %     r_{n-2} - r_{n-1} \cdot q_n &= r_n
    % \end{align*}
    % Substituting gives: for $n \geq 2$,
    % \begin{align*}
    %     r_0 &= a - b \cdot q_0 \\
    %     r_1 &= b - r_0 \cdot q_1 = b - (a - b \cdot q_0) \cdot q_1 = a \cdot (-q_1) + b(1 + q_0 \cdot q_1) \\
    %     r_2 &= r_0 - r_1 \cdot q_{2} = a - b \cdot q_0 - (a \cdot (-q_1) + b(1 + q_0 \cdot q_1)) \cdot q_2
    %     = a \cdot(1 + q_1 \cdot q_2) - b \cdot (q_0 + )
    % \end{align*}
% \end{proof}

\addparttoc{Linear Algebra}
%Studying Finite Dimension Vector Space and Linear map.

\addparttoc{Topology}
\chapter{General Topology}
\section{Topology}
\section{Basis}
\section{Subspace}
\section{Product Space}
\section{Topological Map}
\section{Separation Axiom}
\section{Compactness}
\section{Connectedness}
\section{Tietze Extension Theorem}

\chapter{Metric Space}

\section{Metric Space}

\section{Compact Metric Space}

\section{Complete Metric Space}

\section{Uryson Metrization Theorem}



\chapter{Geometric Topology}
\section{Quotient Space}
\section{Manifold}

\chapter{Algebraic Topology}
\section{Homotopy}
\section{Fundamental Group}

\addparttoc{Analysis}

\chapter{Basic Analysis}

\section{Real Numbers}
\begin{definition}
    Let $\mathcal{R} \defequal \{ \{a_n\} \mid \text{$\{a_n\}$ is Sequence of Rational numbers} \}$. \\
    Define an Equivalent Relation on $\mathcal{R}$ as: for all $\{a_n\}, \{b_n\} \in \mathcal{R}$,
    \begin{center}
        $\{a_n\} \sim \{b_n\}$ \textit{if and only if} for any $\epsilon \in \mathbb{Q}^+$, there exists $N \in \mathbb{N}$ such that $n \geq N \implies |a_n - b_n| < \epsilon$.
    \end{center}
\end{definition}

\section{Basic Topology on $\mathbb{R}$}
\subsection{Representation Theorem for open sets on $\mathbb{R}$}
\begin{definition}
    Let $S$ be an open set of $\mathbb{R}$.\\
    An open interval $I$ is called a \textit{component interval} of $S$ if: $I \subseteq S$ and
    \begin{center}
        There exists no open interval $J$ such that $I \subsetneq J \subseteq S$
    \end{center}
    Note: Component interval is not unique; $S = (0,1) \cup (1,2)$ has two component intervals.
\end{definition}
\begin{lemma}
    
\end{lemma}

\section{Sequence}

\section{Series}

\section{Differentiation}

\section{Riemann Integral}

\section{Sequence of Functions}


\chapter{Measure Theory}
\section{Sigma Algebra}
\section{Measurable}
\section{Lebesgue Integral}

\chapter{Complex Analysis}
\section{Preliminary}
\subsection{Algebra on Complex Numbers}
\subsection{Inequality}
\subsection{Geomeotric Representation}

\section{Analytic Funtions}

\section{Elementary Functions}

\section{Complex Integral}



\chapter{Fourier Analysis}


\chapter{Analysis on Manifolds}

\section{Basic Definition}

\begin{definition}
    Let $A \subset \mathbb{R}^m$; $f:A \to \mathbb{R}^n$, suppose $A$ contains an Open Neiborhood of a point $\mathbf{a}$. \\
    Given non-zero vector $\mathbf{u} \in \mathbb{R}^m$, define \textit{directional derivative} of $f$ at $\mathbf{a}$ with respect to $\mathbf{u}$:
    \begin{align*}
        f'(\mathbf{a}; \mathbf{u}) \defequal \lim_{t \to 0} \frac{f(\mathbf{a} + t \mathbf{u}) - f(\mathbf{a})}{t}
    \end{align*}
    if the limit exists.
\end{definition}

\begin{definition}
    Let $A \subset \mathbb{R}^m$; $f:A \to \mathbb{R}^n$, suppose $A$ contains an Open Neiborhood of a point $\mathbf{a}$. \\
    A map $f$ is said to be \textit{differentiable} at $\mathbf{a}$ if: there exists a linear map $T \in \mathcal{L}(\mathbb{R}^m, \mathbb{R}^n)$ such that
    \begin{align*}
        \frac{f(\mathbf{a} + \mathbf{h}) - f(\mathbf{a}) - T(\mathbf{h})}{\|\mathbf{h}\|} \to \mathbf{0} \as \mathbf{h} \to \mathbf{0}
    \end{align*}
    More rigorously,
    \begin{align*}
        \forall \varepsilon > 0,\
        \exists \delta >0 \st
        \forall \mathbf{h} \in \mathbb{R}^m,\
        0 < \|\mathbf{h}\| < \delta \implies \frac{\|f(\mathbf{a} + \mathbf{h}) - f(\mathbf{a}) - T(\mathbf{h})\|}{\|\mathbf{h}\|} < \varepsilon
    \end{align*}
    The linear map $T$ satisfying the above condition is called \textit{derivative} of $f$ at $\mathbf{a}$.
\end{definition}
\begin{lemma}
    Let $A \subset \mathbb{R}^m$; $f:A \to \mathbb{R}^n$, suppose $A$ contains an Open Neiborhood of a point $\mathbf{a}$. \\
    If linear maps $T, U \in \mathcal{L}(\mathbb{R}^m, \mathbb{R}^n)$ are derivative of $f$ at $\mathbf{a}$, then $T = U$.
\end{lemma}
\begin{proof}
    Suppose $T, U$ are derivative of $f$ at $\mathbf{a}$. \\
    For any $\varepsilon > 0$, there exist $\delta_1, \delta_2 > 0$ such that
    \begin{align*}
        \forall \mathbf{h} \in \mathbb{R}^m,\
        0 < \|\mathbf{h}\| < \delta
        \implies 
        \frac{\|f(\mathbf{a} + \mathbf{h}) - f(\mathbf{a}) - T(\mathbf{h})\|}{\|\mathbf{h}\|},\
        \frac{\|f(\mathbf{a} + \mathbf{h}) - f(\mathbf{a}) - U(\mathbf{h})\|}{\|\mathbf{h}\|} < \frac{\varepsilon}{2}
    \end{align*}
    The Triangle inequality gives that: for small enough $\mathbf{h} \in \mathbf{h}$,
    \begin{align*}
        \frac{\|T(\mathbf{h}) - U(\mathbf{h})\|}{\|\mathbf{h}\|}
        \leq 
        \frac{\|f(\mathbf{a} + \mathbf{h}) - f(\mathbf{a}) - T(\mathbf{h})\|}{\|\mathbf{h}\|} +
        \frac{\|f(\mathbf{a} + \mathbf{h}) - f(\mathbf{a}) - U(\mathbf{h})\|}{\|\mathbf{h}\|}
        < \varepsilon
    \end{align*}
    Now, let $\varepsilon > 0$ be given, and Put $\mathbf{u} \in \mathbb{R}^m$ be a fixed non-zero vector. Then, for small enough $t > 0$,
    \begin{align*}
        \frac{\|T(t \cdot \mathbf{u}) - U(t \cdot \mathbf{u})\|}{\|t \cdot \mathbf{u}\|}
        = \frac{|t| \|T(\mathbf{u}) - U(\mathbf{u})\|}{|t| \|\mathbf{u}\|}
        = \frac{\|T(\mathbf{u}) - U(\mathbf{u})\|}{\|\mathbf{u}\|}
        < \frac{\varepsilon}{\|\mathbf{u}\|}
    \end{align*}
    In sum, For any $\varepsilon > 0$, and for any non-zero vector $\mathbf{u}$, $\|T(\mathbf{u}) - U(\mathbf{u})\| < \varepsilon$.
    Thus $T = U$.
\end{proof}

\newpage
\begin{theorem}
    Let $f:A \to \mathbb{R}^n$ where $A \subseteq \mathbb{R}^m$. \\
    If $f$ is differentiable at $\mathbf{a} \in \mathbb{R}^m$, then all the directional derivative of $f$ at $\mathbf{a}$ exists and
    \begin{align*}
        f'(\mathbf{a} ; \mathbf{u}) = D f(\mathbf{a}) \cdot \mathbf{u}
    \end{align*}
\end{theorem}
\begin{proof}
    Let non-zero $\mathbf{u} \in \mathbb{R}^m$ be given. For any $\varepsilon > 0$, there exists $\delta > 0$ such that
    \begin{align*}
        0 < |t| < \delta \implies \frac{\|f(\mathbf{a} + t \mathbf{u}) - f(\mathbf{a}) - D f(\mathbf{a}) \cdot t \mathbf{u}\|}{\|t \mathbf{u}\|} < \frac{\varepsilon}{\|\mathbf{u}\|}
    \end{align*}
    Arrange above formular,
    \begin{align*}
        \frac{\|f(\mathbf{a} + t \mathbf{u}) - f(\mathbf{a}) - D f(\mathbf{a}) \cdot t \mathbf{u}\|}{\|t \mathbf{u}\|}
        = \frac{1}{|t| \|\mathbf{u}\|} \cdot
        \|f(\mathbf{a} + t \mathbf{u}) - f(\mathbf{a}) - t Df(\mathbf{a}) \cdot \mathbf{u} \|
        = \frac{1}{\|u\|} \cdot
        \left\|\frac{f(\mathbf{a} + t \mathbf{u}) - f(\mathbf{a})}{t} - Df(\mathbf{a}) \cdot \mathbf{u} \right\|
    \end{align*}
    That is, above formular can rearrange to
    \begin{align*}
        0 < |t| < \delta \implies \left\|\frac{f(\mathbf{a} + t \mathbf{u}) - f(\mathbf{a})}{t} - Df(\mathbf{a}) \cdot \mathbf{u} \right\| < \varepsilon
    \end{align*}
\end{proof}

\begin{theorem}
    Let $f:A \to \mathbb{R}^n$ where $A \subseteq \mathbb{R}^m$. \\
    If $f$ is differentiable at $\mathbf{a} \in A$, then $f$ is continuous at $\mathbf{a}$.
\end{theorem}
\begin{proof}
    Denote $B$ be a Derivative of $f$ at $\mathbf{a}$. 
    By Triangle inequality, for any $\mathbf{h} \in \mathbb{R}^m$, 
    \begin{align*}
        \|f(\mathbf{a} + \mathbf{h}) - f(\mathbf{a})\| \leq
        \|f(\mathbf{a} + \mathbf{h}) - f(\mathbf{a}) - B \cdot \mathbf{h}\| + \|B \cdot \mathbf{h}\|
        = \|\mathbf{h}\| \cdot \frac{\|f(\mathbf{a} + \mathbf{h}) - f(\mathbf{a}) - B \cdot \mathbf{h}\|}{\|\mathbf{h}\|} + \|B \cdot \mathbf{h}\|
    \end{align*}
    And the Right side converges to zero, so left side too.
\end{proof}

\begin{definition}
    Let $f:A \to \mathbb{R}$ where $A \subseteq \mathbb{R}^m$. \\
    For an integer $1 \leq j \leq m$, define \textit{$j^{\text{th}}$ Derivative} at $\mathbf{a} \in \mathbb{R}^m$:
    \begin{align*}
        D_j f(\mathbf{a}) \defequal \lim_{t \to 0} \frac{f(\mathbf{a} + t \cdot \mathbf{e_j}) - f(\mathbf{a})}{t}
    \end{align*}
\end{definition}

\begin{theorem}
    Let $f:A \to \mathbb{R}$ where $A \subseteq \mathbb{R}^m$. \\
    If $f$ is differentiable at $\mathbf{a} \in A$, then the Derivative $Df(\mathbf{a})$ can represent as:
    \begin{align*}
        Df(\mathbf{a}) = \begin{bmatrix}
            D_1 f(\mathbf{a}) &
            D_2 f(\mathbf{a}) &
            \cdots
            D_m f(\mathbf{a})
        \end{bmatrix}
    \end{align*}
\end{theorem}
\begin{proof}
    Put the Derivative of $f$ at $\mathbf{a}$:
    \begin{align*}
        Df(\mathbf{a}) = \begin{bmatrix}
            \lambda_1 & \lambda_2 & \cdots & \lambda_m
        \end{bmatrix}
        \in \mathcal{M}_{1, m}
    \end{align*}
    Applying this to above theorem: for each $1 \leq j \leq m$,
    \begin{align*}
        D_j f(\mathbf{a})
        = f'(\mathbf{a}, \mathbf{e}_j)
        = Df(\mathbf{a}) \cdot \mathbf{e}_j
        = \lambda_j
    \end{align*}
    Thus,
    \begin{align*}
        Df(\mathbf{a})
        =
        \begin{bmatrix}
            \lambda_1 & \lambda_2 & \cdots & \lambda_m
        \end{bmatrix}
        =
        \begin{bmatrix}
            D_1 f(\mathbf{a}) & D_2 f(\mathbf{a}) & \cdots & D_m f(\mathbf{a})
        \end{bmatrix}
    \end{align*}
\end{proof}


\newpage
\begin{theorem}
    Let $f: A \to \mathbb{R}^n$ where $A \subseteq \mathbb{R}^m$. For each $j$, let $f_i:A \to \mathbb{R}: \mathbf{x} \mapsto p_i(\mathbf{x})$ so that
    \begin{align*}
        f(\mathbf{x}) = 
        \begin{bmatrix}
            f_1(\mathbf{x}) \\
            f_2(\mathbf{x}) \\
            \vdots \\
            f_n(\mathbf{x})
        \end{bmatrix}
    \end{align*}
    Then, following equivalent statement be true:
    \begin{center}
        $f$ is differentiable at $\mathbf{a}$ \textit{if and only if}
        all $f_i$ is differentiable at $\mathbf{a}$
    \end{center}
    Moreover, if $f$ is differentiable at $\mathbf{a}$,
    \begin{align*}
        Df(\mathbf{a})
        = \begin{bmatrix}
            Df_1(\mathbf{x}) \\
            Df_2(\mathbf{x}) \\
            \vdots \\
            Df_n(\mathbf{x})
        \end{bmatrix}
    \end{align*}
\end{theorem}

\addparttoc{Differential Geometry}



\addparttoc{Differential Equation}



\addparttoc{}








\end{document}


%cd /Users/persist/Desktop/My_Mathematics_Book
%latexmk -pdf -bibtex main.tex
%latexmk -xelatex -synctex=1 -interaction=nonstopmode main.tex